\documentclass[final]{siamart0516}
% SIAM class siamart0516 (May 2016 macros; obtain the full distribution from SIAM).
% The class loads amsmath, ntheorem, hyperref, xcolor and cleveref itself.
\usepackage{amssymb,mathtools}
\usepackage{enumitem}
\usepackage{booktabs,longtable,tabularx,array}
\usepackage{tikz}
\usepackage{adjustbox}
\usetikzlibrary{positioning,arrows.meta,shapes.misc}

\setlist{itemsep=2pt,topsep=4pt}
\setlength{\emergencystretch}{2.5em}
\newcolumntype{L}[1]{>{\raggedright\arraybackslash}p{#1\linewidth}}
\setlength{\LTpre}{6pt}\setlength{\LTpost}{6pt}

% ------------------------------------------------------------------ theorem environments
\numberwithin{theorem}{section}
\numberwithin{equation}{section}
\newsiamthm{conjecture}{Conjecture}
\newsiamthm{claim}{Claim}
\newsiamremark{remark}{Remark}
\newsiamremark{example}{Example}
\newsiamremark{problem}{Open Problem}
% Unnumbered remark* and definition* (ntheorem's own starred variants are numbered by this class).
\makeatletter
\newtheoremstyle{siamnonum}%
 {\item[\hskip\labelsep\hskip\parindent\theorem@headerfont ##1\theorem@separator]}%
 {\item[\hskip\labelsep\hskip\parindent\theorem@headerfont ##1\ \textup{(##3)}\theorem@separator]}
\makeatother
\theoremstyle{siamnonum}
\theoremheaderfont{\normalfont\itshape}
\theorembodyfont{\normalfont}
\theoremseparator{.}
\theoremsymbol{}
\newtheorem{remarkstar}{Remark}
\theoremheaderfont{\normalfont\sc}
\theorembodyfont{\normalfont\itshape}
\newtheorem{definitionstar}{Definition}
\renewenvironment{remark*}{\remarkstar}{\endremarkstar}
\renewenvironment{definition*}{\definitionstar}{\enddefinitionstar}
\theoremstyle{plain}
\crefname{conjecture}{Conjecture}{Conjectures}
\crefname{claim}{Claim}{Claims}
\crefname{problem}{Open Problem}{Open Problems}
\crefname{remark}{Remark}{Remarks}
\crefname{example}{Example}{Examples}
\providecommand{\qed}{\hfill\proofbox}

% ------------------------------------------------------------------ macros
\newcommand{\code}[1]{{\ttfamily\nolinkurl{#1}}}
\makeatletter\g@addto@macro\UrlBreaks{\do\_\do\-}\makeatother
\newcommand{\NN}{\mathbb{N}}
\newcommand{\ZZ}{\mathbb{Z}}
\newcommand{\QQ}{\mathbb{Q}}
\newcommand{\RR}{\mathbb{R}}
\newcommand{\CC}{\mathbb{C}}
\newcommand{\FF}{\mathbb{F}}
\newcommand{\Qz}{\QQ(z)}
\newcommand{\Rz}{\RR(z)}
\newcommand{\Qzbar}{\overline{\QQ(z)}}
\newcommand{\Puis}{\mathcal{P}_{\RR}}
\newcommand{\val}{\mathcal{V}}            % the value
\newcommand{\pay}{\mathcal{U}}            % the payoff
\newcommand{\SG}{\mathcal{G}}             % the state graph
\newcommand{\arena}{\mathcal{A}}          % an arena
\newcommand{\Mpot}{\mathcal{M}}           % set-potential values
\newcommand{\cont}{\mathcal{K}}           % continuants
\newcommand{\game}{\mathfrak{G}}          % the game G(H;A)
\newcommand{\arcs}{\mathsf{A}}            % a set of one-way arcs
\newcommand{\mm}{\mathfrak{m}}            % number of maximum matchings
\newcommand{\Dv}{\mathsf{D}}              % divisor map
\newcommand{\ord}{\operatorname{ord}}
\newcommand{\Res}{\operatorname{Res}}
\newcommand{\back}{\operatorname{back}}
\newcommand{\tr}{\operatorname{tr}}
\renewcommand{\Im}{\operatorname{Im}}
\newcommand{\rk}{\operatorname{rk}}
\newcommand{\poly}{\operatorname{poly}}
\newcommand{\Sym}{\operatorname{Sym}}
\newcommand{\Rdown}{R^{\downarrow}}
\newcommand{\nplus}{N^{+}}
\newcommand{\cls}[1]{\mathsf{#1}}
\newcommand{\PSPACE}{\cls{PSPACE}}
\newcommand{\EXP}{\cls{EXP}}
\newcommand{\EXPTIME}{\cls{EXPTIME}}
\newcommand{\NP}{\cls{NP}}
\newcommand{\coNP}{\cls{coNP}}
\newcommand{\BPP}{\cls{BPP}}
\newcommand{\BQP}{\cls{BQP}}
\newcommand{\Pclass}{\cls{P}}
\newcommand{\NC}{\cls{NC}}
\newcommand{\CH}{\cls{CH}}
\newcommand{\MA}{\cls{MA}}
\newcommand{\IP}{\cls{IP}}
\newcommand{\sharpP}{\#\cls{P}}
\newcommand{\Nw}{\mathrm{N}}   % N-position
\newcommand{\Pw}{\mathrm{P}}   % P-position
\newcommand{\defeq}{\coloneqq}
\newcommand{\one}{\mathbf{1}}
